%% ==========================================================================
%% DRAFT -- August 12, 2026. Author byline: Daniel Kirtchakov (05oz,
%% halfounce.io; no institutional affiliation; ORCID 0009-0009-5213-4098).
%%
%% Framing constraints honoured (PROTOCOL section 11): the note opens with the
%% mathematical question, states the certified object as real mathematics with
%% the Newell formulas and the interval enclosure made precise, reports the
%% pathology quantification with rigorous numbers, separates what is certified
%% from what is enclosed and what is trusted, cites the computation in a single
%% paragraph, and closes with the questions it opens. Every reported digit-loss
%% figure is a rigorous consequence of the shipped two-sided enclosure.
%% ==========================================================================
\documentclass[11pt]{amsart}

\usepackage[margin=1in]{geometry}
\usepackage{amsmath,amssymb}
\usepackage{booktabs}
\usepackage{array}
\usepackage[colorlinks=true,allcolors=blue]{hyperref}

\emergencystretch=3em
\tolerance=1200
\hbadness=4000
\addtolength{\skip\footins}{0pt minus 3pt}

\newcommand{\Nxx}{N_{xx}}
\newcommand{\Nxy}{N_{xy}}
\newcommand{\Rnorm}{R}
\newcommand{\asinh}{\operatorname{asinh}}
\DeclareMathOperator{\atan}{atan}

\theoremstyle{plain}
\newtheorem{theorem}{Theorem}[section]
\newtheorem{lemma}[theorem]{Lemma}
\newtheorem{proposition}[theorem]{Proposition}
\theoremstyle{definition}
\newtheorem{definition}[theorem]{Definition}
\theoremstyle{remark}
\newtheorem{remark}[theorem]{Remark}
\newtheorem{question}[theorem]{Question}

\begin{document}

\title[Certified Newell demagnetization-tensor tables]{Certified
demagnetization-tensor reference tables: two-sided rational enclosures of the
Newell entries, and a rigorous map of where double-precision micromagnetics
loses its digits}

\author{Daniel Kirtchakov}
\address{Independent researcher (\texttt{05oz}), Half Ounce Research; no
institutional affiliation. ORCID 0009-0009-5213-4098.
\texttt{daniel@halfounce.io} $\cdot$ \texttt{https://halfounce.io}}

\date{Draft of August 12, 2026; revised November 10, 2026 (v0.17.0).}

\begin{abstract}
OOMMF and other finite-difference micromagnetic simulators build their
demagnetizing field from the Newell demagnetization tensor of a pair of uniformly
magnetized rectangular cells; some evaluate Newell's closed form, while MuMax3
computes the tensor by numerical integration. It is a textbook fact, documented but
uncertified, that the closed-form evaluation of this tensor loses all
significant digits by catastrophic cancellation once the two cells are more
than a few hundred cell widths apart. We ask whether the tensor entries can be
pinned rigorously --- by two-sided rational enclosures a skeptic re-derives from
a small certificate using nothing but a standard-library program --- across the
whole regime a real simulation uses, including the large-separation regime
where double precision fails. We give such a certificate. Each entry carries an
interval $[N_{\mathrm{lo}},N_{\mathrm{hi}}]$ of dyadic rationals with
$N_{\mathrm{lo}}\le N_{\mathrm{true}}\le N_{\mathrm{hi}}$, obtained by
evaluating Newell's formulas in interval arithmetic that rounds outward on every
operation; the transcendental pieces (\texttt{sqrt}, \texttt{atan},
\texttt{log}) are enclosed by rigorously-truncated series, and everything else
is exact rational arithmetic. Against these enclosures we measure, rigorously,
the two floating-point routes a simulator can take: the naive double-precision
analytic evaluation and a truncated asymptotic expansion. The naive analytic route
degrades by about six correct decimal digits per decade of separation --- from
$\sim\!15$ digits at one cell to \emph{zero} correct digits near $300$ cells and
a value with the wrong sign and wrong order of magnitude beyond --- exactly the
$r^{6}$ law and the $\sim\!300$-cell breakdown that the accuracy literature
predicts but never certified; the asymptotic route is complementary, poor at
short range and good at long range, and the crossover is located rigorously.
The enclosures remain tight to tens of digits everywhere, including where double
precision has none. A certificate JSON and a standard-library checker are the
public unit; the checker re-derives every enclosure, the naive value and every
digit-loss figure from scratch, and rejects tampering both at the hash layer and
at the mathematical re-derivation layer. We state what is certified, what
is merely enclosed, and what is trusted, and close with the questions the table
opens. A November 2026 addendum (\S\ref{sec:bound}) proves an explicit bound on
the truncation error of the asymptotic series, valid for every rectangular cell,
every tensor component and every separation larger than the cell diagonal.
\emph{(Correction, v0.17.0, 2026-11-10: through v0.16.0 the abstract called the
second route ``OOMMF's asymptotic expansion''. The values measured are for a
three-term truncation of the series; OOMMF carries a fourth term. See
\S\ref{sec:results}. Its first sentence also counted MuMax3, which integrates
numerically, among codes using Newell's closed form.)}
\end{abstract}

\maketitle

{\small
\noindent\textbf{Computation and authorship.} All formula transcription,
interval-arithmetic implementation, verification, and drafting in this work were
produced by Claude (Anthropic), directed by the author, on a single Apple laptop.
The public artifact is a certificate together with a checker that imports only
the Python standard library and, though not code-independent of the private
generator (see \S\ref{sec:trusted}), recomputes every enclosure, the naive double
value, and every digit-loss figure from the certificate alone.

\smallskip
\noindent\textbf{Prior-art record.} The primary sources (Newell, Williams and
Dunlop 1993; OOMMF's \texttt{demagcoef.cc}; Chernyshenko and Fangohr 2015;
Bj\o rk and d'Aquino 2023) were read at the relevant sections on August 11--12,
2026; the passages quoted are verbatim. Novelty was swept the same day against
arXiv and the web; the dated record ships with the artifacts.
\emph{(Correction, v0.17.0, 2026-11-10: the asymptotic section of
\texttt{demagcoef.cc} was not among the sections read. The asymptotic form used
in this note was transcribed from the demagderiv notebook \cite{demagderiv},
which follows a 2010 pre-release of OOMMF; see \S\ref{sec:results}. The November
addendum has its own dated record,
\texttt{SWEEP-RECORD-DEMAG-BOUND-2026-11-10.md}.)}\par}

\section{The question}

A finite-difference micromagnetic simulation spends almost all of its arithmetic
on one thing: the demagnetizing (magnetostatic) field. In the energy-based
discretization used by OOMMF and its descendants the field is a discrete
convolution of the magnetization with a fixed kernel, the \emph{demagnetization
tensor} $\mathbf N(\mathbf r)$, whose entry $N_{\alpha\beta}(\mathbf r)$
measures the magnetostatic interaction between two uniformly magnetized
rectangular cells separated by $\mathbf r$. The kernel is computed once, at
setup, from Newell's closed-form analytic formulas \cite{Newell1993}, and it is
then reused for the life of the simulation. Every reported number a
micromagnetic study produces --- a switching field, a coercivity, a resonance
frequency, an energy barrier --- inherits whatever error is in that kernel.

The kernel has a known pathology. Newell's formulas express each tensor entry as
a fixed linear combination of a transcendental ``governing function'' $f$ (for
the diagonal entries) or $g$ (for the off-diagonal entries) evaluated at the
corners of the two cells. Each governing-function value grows like $r^{3}$ with
the separation $r$, while their combination --- the tensor entry itself --- decays
like $1/r^{3}$. Evaluated in floating point, the combination is a difference of
large nearly-equal quantities, and the cancellation is catastrophic.
Chernyshenko and Fangohr state the consequence precisely: the relative error of
the double-precision analytic formula ``is of the order $10^{-15}r^{6}$,'' so
``for cell separations greater than $10^{15/6}\approx 300$ the analytical
computation will contain no significant digits at all'' \cite[\S II.B]{CF2015}.
The simulation packages know this and route around it: OOMMF switches to
asymptotic expansions beyond a switching radius \cite{CF2015}. Since version
2.0a2 it computes a schedule of such radii from the cell dimensions and evaluates
the analytic formula in double-double ($106$-bit) arithmetic below the innermost
one \cite{OOMMF}; earlier releases take the analytic differences in double or
long double, depending on the release and the platform. That routing is itself
uncertified, and the codes that do neither inherit the failure directly.
\emph{(Correction, v0.17.0, 2026-11-10: through v0.16.0 this sentence did not
date the double-double evaluation.)}

This note asks a different question. Not ``what is the fastest accurate
kernel?'' but: \emph{can each tensor entry be bracketed rigorously, by a rational
interval a stranger re-derives from a small certificate using only a
standard-library program, across the entire regime a simulation uses --- and can
the floating-point failure then be measured, rather than merely predicted?} The
cancellation that defeats double precision is not an obstruction to a
certificate; it is a statement about working precision, and it is defeated
outright by carrying the same arithmetic in exact rationals with outward-rounded
transcendental enclosures. The output is, for each cell aspect ratio, offset
direction, separation, and tensor component, a pair of dyadic rationals
$N_{\mathrm{lo}}\le N_{\mathrm{hi}}$ with
\[
   N_{\mathrm{lo}} \;\le\; N_{\alpha\beta}(\mathbf r) \;\le\; N_{\mathrm{hi}},
\]
together with the naive double value at the same point and a rigorous bracket,
derived from the enclosure, on how many of its digits are correct. The
comparison is the contribution: it tells a simulation author exactly where, and
by how much, their tensor goes wrong.

\section{The certified object}\label{sec:object}

\paragraph{The tensor.}
Following the convention of \cite{CF2015}, the demagnetization tensor of two
cells of common volume $|\sigma|=\Delta x\,\Delta y\,\Delta z$ with centre
offset $\mathbf r$ is
\begin{equation}
  \mathbf N(\mathbf r) \;=\; \frac{1}{4\pi|\sigma|}
  \int_{\sigma_1}\!\!\int_{\sigma_2(\mathbf r)}
  \nabla_{\mathbf r_1}\nabla_{\mathbf r_2}\frac{1}{|\mathbf r_1-\mathbf r_2|}
  \,\mathrm d\mathbf r_2\,\mathrm d\mathbf r_1 ,
  \label{eq:tensordef}
\end{equation}
so that the demagnetizing field is $\mathbf H=-\mathbf N\cdot\mathbf M$; in
Newell's own notation \cite{Newell1993} this is $H_x=-N_{xx}M_x$, his
formula~(16). The tensor is symmetric and trace-free for $\mathbf r\neq 0$, and
the diagonal self-term $\mathbf r=0$ obeys $N_{xx}+N_{yy}+N_{zz}=1$.
\emph{(Correction, v0.17.0, 2026-11-10: through v0.16.0 equation~\eqref{eq:tensordef}
carried a leading minus sign, copied from \cite[eq.~(2)]{CF2015}. With one
gradient taken in each integration variable that sign gives the opposite of the
values computed here and in OOMMF, for which two cells in line along $x$ have
$N_{xx}<0$. The certified values were always computed from Newell's closed form
with the correct sign. The \texttt{formulas.tensor\_def} text of the released
\texttt{demag\_certificate.json} carries the same superseded sign and is left
unchanged so that the record stays byte-identical.)}

\paragraph{Newell's governing functions.}
The double volume integral in \eqref{eq:tensordef} has a closed form. For the
diagonal entry, with $R=\sqrt{x^2+y^2+z^2}$,
\begin{equation}
  f(x,y,z) = \tfrac16(2x^2-y^2-z^2)R
  + \tfrac{y}{2}(z^2-x^2)\,\asinh\!\frac{y}{\sqrt{x^2+z^2}}
  + \tfrac{z}{2}(y^2-x^2)\,\asinh\!\frac{z}{\sqrt{x^2+y^2}}
  - xyz\,\atan\!\frac{yz}{xR},
  \label{eq:newellf}
\end{equation}
and for the off-diagonal entry,
\begin{align}
  g(x,y,z) = {}&{-}\tfrac13 xyR
  - \tfrac{z^3}{6}\atan\!\frac{xy}{zR}
  - \tfrac{z y^2}{2}\atan\!\frac{xz}{yR}
  - \tfrac{z x^2}{2}\atan\!\frac{yz}{xR}
  + xyz\,\asinh\!\frac{z}{\sqrt{x^2+y^2}} \notag\\
  &{}+ \tfrac{y}{6}(3z^2-y^2)\asinh\!\frac{x}{\sqrt{y^2+z^2}}
  + \tfrac{x}{6}(3z^2-x^2)\asinh\!\frac{y}{\sqrt{x^2+z^2}} ,
  \label{eq:newellg}
\end{align}
where $\asinh t=\ln\!\big(t+\sqrt{1+t^2}\big)$. These are the functions
\texttt{Oxs\_Newell\_f} and \texttt{Oxs\_Newell\_g} of OOMMF's
\texttt{demagcoef.cc}; we implement them in the mathematically equivalent
logarithm form the code uses, $\asinh\!\frac{y}{\sqrt{x^2+z^2}}
=\tfrac12\ln\!\frac{(y+R)^2}{x^2+z^2}$, and we mirror every guard the code takes
when an argument coordinate vanishes.

\paragraph{The stencil.}
The tensor entry is a second difference of the governing function in each of the
three axes, at step equal to the cell edge:
\begin{equation}
  \Nxx(\mathbf r) = -\frac{1}{4\pi\,\Delta x\,\Delta y\,\Delta z}\,
  \delta^{2}_{x}\delta^{2}_{y}\delta^{2}_{z}\,[\,f\,](\mathbf r),
  \qquad
  \Nxy(\mathbf r) = -\frac{1}{4\pi\,\Delta x\,\Delta y\,\Delta z}\,
  \delta^{2}_{x}\delta^{2}_{y}\delta^{2}_{z}\,[\,g\,](\mathbf r),
  \label{eq:stencil}
\end{equation}
where $\delta^{2}_{x}[f](x,y,z)=f(x+\Delta x,y,z)-2f(x,y,z)+f(x-\Delta x,y,z)$, and
$\delta^{2}_{y}$, $\delta^{2}_{z}$ likewise. With its leading minus sign,
\eqref{eq:stencil} is a $27$-point stencil with weights $-1$ at the eight
corners, $+2$ at the twelve edges, $-4$ at the six faces and $+8$ at the centre,
as in OOMMF's \texttt{Oxs\_CalculateNxx}. The remaining diagonal and
off-diagonal entries are coordinate relabelings:
$N_{yy}(\mathbf r;\Delta)=N_{xx}$ with $(x,\Delta x)\leftrightarrow(y,\Delta y)$,
$N_{zz}$ with $x\leftrightarrow z$, and $N_{xz},N_{yz}$ from $N_{xy}$ likewise.
The self-term $\mathbf r=0$ uses Newell's separate closed form for $N_{xx}(0)$
(\texttt{Oxs\_SelfDemagNx}), which we also carry. \emph{(Correction, v0.17.0,
2026-11-10: through v0.16.0 equation~\eqref{eq:stencil} had no leading minus sign
and did not define the second difference. Read with the standard second
difference above, it gave the opposite of the weights stated with it. The
weights, which the certified values use, are unchanged; the
\texttt{formulas.stencil} text of the released \texttt{demag\_certificate.json}
carries the old form and is left byte-identical.)}

\paragraph{The enclosure, and what is certified.}
We evaluate \eqref{eq:newellf}--\eqref{eq:stencil} in interval arithmetic over
dyadic-rational endpoints. Every interval $[a,b]$ rigorously encloses the true
value; every operation --- addition, multiplication, division, and the three
transcendentals --- rounds its result \emph{outward} to the grid of denominator
$Q=2^{P}$ (working precision $P=256$ bits, about $77$ decimal digits), which
keeps the rationals bounded and makes the whole computation deterministic. The
transcendentals are enclosed from below and above:
\begin{itemize}
  \item $\sqrt q$ for a nonnegative rational $q$ by integer square root at the
  grid scale: with $m=\lfloor\sqrt{\lfloor qQ^2\rfloor}\rfloor$ one has
  $(m/Q)^2\le q<((m{+}1)/Q)^2$, a two-sided rational bound;
  \item $\atan t$ by half-angle reduction $\atan t = 2\atan\!\big(t/(1+\sqrt{1+t^2})\big)$
  until $|t|\le\tfrac12$, then the alternating Maclaurin series, whose truncation
  error is bounded by the first omitted term;
  \item $\ln p$ by pulling out powers of two until the mantissa lies in
  $[\tfrac23,\tfrac43]$, then $2\operatorname{atanh}\frac{m-1}{m+1}$ with a
  geometric-series tail bound; $\ln 2$ and
  $\pi=4\big(4\atan\tfrac15-\atan\tfrac1{239}\big)$ (Machin's formula) are
  certified once at the same precision. \emph{(Correction, v0.17.0, 2026-11-10:
  through v0.16.0 the expression printed here was $\pi/4$; the \texttt{method}
  text of the released \texttt{demag\_certificate.json} carries it unchanged.)}
\end{itemize}
What is \emph{certified} is the two-sided inequality $N_{\mathrm{lo}}\le
N_{\alpha\beta}\le N_{\mathrm{hi}}$ for the analytic Newell tensor entry. What is
\emph{enclosed} rather than exact are exactly the transcendental pieces, each to
rigorous rational bounds; the stencil weights, the cell dimensions, the offsets
and all the surrounding arithmetic are exact. As a numerical anchor,
\texttt{demagcoef.cc} carries in its header a table of $N_{xx}$ and $N_{xy}$
values Donahue computed to $50$ digits in Maple --- $31$ rows, $24$ of them
nonzero --- and the $16$ of those rows that form our anchor agree with the
midpoints of our enclosures to at least $49.6$ digits; containment would be the
wrong test,
since the enclosures are roughly $77$ digits tight and a $50$-digit value
therefore lies outside them in general, so agreement at the gold values' own
precision is what is checked. An independent
$220$-digit floating-point evaluation of the same formulas does lie inside every
enclosure of the table.\footnote{Erratum, v0.12.1 (2026-08-12): as first
released in v0.12.0 this sentence asserted that the $16$ Maple values
\emph{lie inside} the enclosures. That is false for the reason just given —
$15$ of the $16$ lie outside, the exact zero being the exception — and the
build log (\texttt{FIXLOG.md}, decision S4) had already recorded the correct
test. The agreement figure of $49.6$ digits is unaffected, as is every
certified enclosure. Only this prose claim changed. Erratum, 2026-08-12: through v0.12.1 the same sentence read ``the $16$ diagonal and
off-diagonal values Donahue computed to $50$ digits in Maple and recorded in
\texttt{demagcoef.cc},'' which misdescribes the source table --- the header
carries $31$ rows, $24$ of them nonzero, and $16$ is the size of the subset we
anchor against, not of the table. The anchor set, the agreement figure and the
enclosures are unchanged.}

\begin{remark}[Internal consistency, also certified]\label{rem:consistency}
The enclosures satisfy the tensor's own structural identities, and the checker
can test them without any external reference. At every mutual point the interval
sum $[N_{xx}]+[N_{yy}]+[N_{zz}]$ encloses $0$ (trace-freeness for
$\mathbf r\neq0$); at the self-term the same sum encloses $1$; and the
permutation identities hold as interval containments among the components the
table carries, since $N_{yy},N_{zz},N_{xz}$ are computed by relabeling the
\emph{same} $f$ and $g$ evaluations that produce $N_{xx}$ and $N_{xy}$
($N_{yz}$ is carried nowhere, so no identity involving it is tested). These are not
independent physical facts to be trusted --- they are consequences of
\eqref{eq:newellf}--\eqref{eq:stencil} that a rigorous enclosure must respect, and
their violation would signal a defect in the enclosure itself. Every one holds
across the table.
\end{remark}

\section{Where this sits}

The analytic formulas are Newell, Williams and Dunlop's \cite{Newell1993}; the
reference implementation is Donahue and Porter's OOMMF \cite{OOMMF}, whose
analytic routines in \texttt{demagcoef.cc} we follow line for line, and whose
double-double evaluation, four-term asymptotic form and model-based switching
radii are the state of the practice. \emph{(Correction, v0.17.0, 2026-11-10:
through v0.16.0 this sentence did not limit ``line for line'' to the analytic
routines. The asymptotic form used in this note was not taken from
\texttt{demagcoef.cc}; see \S\ref{sec:results}.)} The large-separation
cancellation is analysed by Chernyshenko and Fangohr \cite{CF2015}, who give the
$10^{-15}r^{6}$ error law and the $\approx 300$-cell breakdown, and propose
numerical integration as an alternative; Bj\o rk and d'Aquino \cite{Bjork2023}
measure the accuracy of the analytic tensor against the dipolar field across
prism, tetrahedron and cylinder tiles and report accuracy \emph{plots} out to
$10^{4}$ tile radii for single and double precision. Both use high-precision
floating point internally to obtain a reference value ``exact up to machine
precision'' for their error estimates. None of this work produces a
\emph{certified} object: the reference values are floating-point, not two-sided
rigorous enclosures; they are used inside an error study, not published as a
table a third party can rely on; and there is no checker and no artifact. Our
same-day novelty sweep (\S\ref{sec:cert}) found no certified or interval-arithmetic
demagnetization-tensor table anywhere. What is new here is the intersection: a
rigorous two-sided enclosure of the Newell entries across the simulation regime,
a rigorous measurement of the double-precision and asymptotic failures against
those enclosures, and re-verification of every entry by a standard-library
checker that is not code-independent of the generator (\S\ref{sec:trusted}). We
claim that intersection, and no new tensor value or physical result within it.
\emph{(Correction, v0.17.0, 2026-11-10: through v0.16.0 this sentence called the
re-verification independent; see \S\ref{sec:trusted}.)}

\section{Results}\label{sec:results}

We instantiate the certificate over four cell aspect ratios --- a cube
$(1,1,1)$, a thin film $(1,1,\tfrac14)$, an elongated cell $(1,1,4)$, and a
rectangular in-plane cell $(2,1,1)$ --- four offset directions (along $x$, along
$z$, the in-plane face diagonal, and the body diagonal), and separations from
one to ten thousand cells, carrying a selected set of tensor components at each
point together with the diagonal self-terms at the origin --- $862$ certified
entries in all: $850$ across fifty geometry-and-component combinations (cell,
offset direction, tensor component) away from the origin, and twelve self-terms,
$N_{xx}(0),N_{yy}(0),N_{zz}(0)$ for each of the four cells. The selection is not
the whole tensor at every point. A component that vanishes identically at a
point is dropped there --- no carried enclosure contains $0$ --- and $N_{yy}$ is
dropped off the coordinate axes, where it either vanishes or is returned by
trace-freeness from the certified $N_{xx}$ and $N_{zz}$; $N_{yz}$ is carried
nowhere in the table. For the cube, the film and the elongated cell
$N_{yz}$ is a relabelling of a component the table does carry, but for the
rectangular cell on the body diagonal it is neither zero nor a relabelling ---
at offset $(4,2,2)$ it is $\approx-8.15\times10^{-4}$ --- so that one component,
at that one cell and direction, lies outside the certified
table.\footnote{Erratum, 2026-08-12: v0.12.0 and v0.12.1 described the
grid as ``keeping the nonzero components at each point.'' That is false as
stated: the table carries no $N_{yz}$ anywhere, and for the rectangular cell on
the body diagonal $N_{yz}$ is nonzero and is not a relabelling of any carried
component. The coverage is restated above, and the count of $862$ entries, every
certified enclosure and every digit-loss bracket are unchanged. Only this prose
claim changed.}
Table~\ref{tab:pathology} reads the central result off the canonical case: the
cube, on-axis $N_{xx}$.

\begin{table}[h]
\centering
\caption{Cube cell, on-axis offset $\mathbf r=(n,0,0)$, component $N_{xx}$. The
certified enclosure midpoint (all endpoints are exact dyadic rationals in the
certificate); the naive IEEE-754 double analytic value; and rigorous brackets on
the number of correct significant digits each floating-point route retains,
derived from the enclosure. ``Naive'' is the double-precision analytic Newell
stencil; ``three terms'' and ``four terms'' are the asymptotic
series~\eqref{eq:series} truncated after $1/r^{7}$ (the value stored in the
certificate) and after $1/r^{9}$ (as in OOMMF's \texttt{demagcoef.cc}, without
its cell subdivision; computed by \texttt{demag-scripts/four\_term\_digits.py}).
From about $100$ cells (four terms) and $300$ cells (three terms) the truncations
sit at the double-precision floor, where the last digit depends on the order of
the floating-point operations. \emph{(Correction, v0.17.0, 2026-11-10: through
v0.16.0 this caption called the three-term column OOMMF's asymptotic expansion;
see \S\ref{sec:results}.)}}
\label{tab:pathology}
\begin{tabular}{r r r r r}
\toprule
$n$ (cells) & certified $N_{xx}$ (midpoint) & naive double & three terms &
four terms \\
 & & \multicolumn{3}{c}{correct digits} \\
\midrule
$1$      & $-1.35017180544\times10^{-1}$ & $15.2$ & $0.5$ & $0.6$ \\
$5$      & $-1.27235750359\times10^{-3}$ & $11.2$ & $5.1$ & $5.7$ \\
$13$     & $-7.24408320619\times10^{-5}$ & $8.6$  & $7.7$ & $9.0$ \\
$21$     & $-1.71854643363\times10^{-5}$ & $7.1$  & $9.0$ & $10.7$ \\
$55$     & $-9.56603669319\times10^{-7}$ & $5.4$  & $11.5$ & $14.0$ \\
$100$    & $-1.59154942396\times10^{-7}$ & $5.1$  & $13.1$ & $16.7$ \\
$300$    & $-5.89462752160\times10^{-9}$ & $\mathbf{0.4}$ & $17.2$ & $17.2$ \\
$500$    & $-1.27323954473\times10^{-9}$ & $-1.0$ & $16.0$ & $16.0$ \\
$1000$   & $-1.59154943092\times10^{-10}$ & $-2.9$ & $16.4$ & $16.4$ \\
$10000$  & $-1.59154943092\times10^{-13}$ & $-8.9$ & $16.3$ & $16.3$ \\
\bottomrule
\end{tabular}
\end{table}

\paragraph{The naive analytic route, measured.}
The correct-digit count of the double-precision analytic evaluation falls by
almost exactly six per decade of separation --- the signature of the
$10^{-15}r^{6}$ relative-error law: from $15.2$ digits at one cell to $-8.9$ at
$n=10^{4}$ is $24.1$ digits over four decades, $6.0$ per decade. It drops below
a single correct digit near $n=300$ ($0.4$ digits, i.e. essentially no correct
significant figure, precisely the breakdown \cite{CF2015} predicts) and below
zero by $n=500$. Beyond that the value is not merely imprecise but wrong: at
$n=1000$ the naive $-1.14\times10^{-7}$ is three orders of magnitude too large,
and at $n=10^{4}$ it is $+1.17\times10^{-4}$ --- wrong sign, nine orders of
magnitude too large --- while the true value is $-1.59\times10^{-13}$. This is a
kernel entry that a from-scratch double-precision analytic implementation returns
as garbage, certified as garbage against a rigorous enclosure.

\paragraph{The breakdown radius is not a single number.}
The ``$\approx 300$ cells'' of \cite{CF2015} is a cube-scale rule of thumb, and
the certificate shows it is exactly that. Across the four aspect ratios and four
directions, the separation at which the naive double first reaches zero correct
digits ranges from about $100$ cells (the elongated cell, off-diagonal
components) to about $2000$ cells (the thin film's normal component $N_{zz}$,
which is large and so tolerates more cancellation before the relative error
reaches one). A thin film loses its \emph{in-plane} $N_{xx}$ near $300$ cells but
holds its \emph{out-of-plane} $N_{zz}$ to nearly $2000$; an elongated cell is the
reverse. The certificate records the crossing separation for each of the $50$
geometry/component combinations, turning a folklore constant into a per-geometry
map; Table~\ref{tab:breakdown} shows a representative slice.

\begin{table}[h]
\centering
\caption{The breakdown radius is geometry-dependent. For a representative set of
cell/component pairs, the smallest separation $n$ (in cells) at which the naive
double analytic value first retains \emph{zero} correct significant digits ---
rigorously determined from the certified enclosure. The ``$\approx 300$'' of
\cite{CF2015} is the cube scale; real geometries span an order of magnitude
around it.}
\label{tab:breakdown}
\begin{tabular}{l l r}
\toprule
cell $(\Delta x,\Delta y,\Delta z)$ & component (direction) & breakdown $n$ \\
\midrule
cube $(1,1,1)$          & $N_{xx}$ on-axis                 & $500$ \\
thin film $(1,1,\tfrac14)$ & $N_{xx}$ in-plane (along $x$) & $300$ \\
thin film $(1,1,\tfrac14)$ & $N_{zz}$ normal (along $z$)   & $2000$ \\
elongated $(1,1,4)$     & $N_{xx}$ (along $x$)             & $1000$ \\
elongated $(1,1,4)$     & $N_{zz}$ (along $z$)             & $200$ \\
rectangular $(2,1,1)$   & $N_{xy}$ (face diagonal $xy$)    & $200$ \\
\bottomrule
\end{tabular}
\end{table}

\paragraph{The asymptotic route, and the crossover.}
A truncated asymptotic expansion is the mirror image of the analytic route. It
is \emph{poor} at short range ($0.5$ correct digits at $n=1$, where the $1/r$
expansion has no right to converge) and good at long range; in
Table~\ref{tab:pathology} it reaches its double-precision floor, near $16$
digits, from about $100$ cells for four terms and about $300$ for three, and for
other cells and directions it can take longer (for the film $(1,1,\tfrac14)$
along its normal, about $500$ cells for four terms and $2{,}000$ for three). The two
floating-point routes therefore cross. For the cube on-axis $N_{xx}$ and the
three-term form, the naive analytic value is the more accurate up to about
thirteen cells and the asymptotic from about twenty-one on ($8.6$ versus $7.7$
digits at $n=13$; $7.1$ versus $9.0$ at $n=21$). The four-term form crosses
earlier: $9.6$ versus $7.3$ digits at $n=8$, $8.6$ versus $9.0$ at $n=13$. Both
are dominated everywhere by the certified enclosure, and because the enclosure
brackets the truth, the crossover radius of a given truncation is itself
bracketed rigorously by the certificate, per aspect ratio and direction. These
crossovers describe a code that evaluates both routes in double precision.
Since version 2.0a2 OOMMF evaluates the analytic formula in double-double below
its switching radii, chooses those radii from assumed error and speed models,
and sums subdivided-cell asymptotics in a middle range; earlier releases take the
differences in double or long double (\S1).
\emph{(Correction, v0.17.0, 2026-11-10: through v0.16.0 this paragraph called the
three-term form ``OOMMF's asymptotic expansion''. It is a three-term truncation
transcribed from the demagderiv notebook \cite{demagderiv}, which follows a 2010
pre-release of OOMMF. Every \texttt{demagcoef.cc} on the public mirror of OOMMF's
releases, from 1.2a5bis to 2.1a0, carries a fourth term, $1/r^{9}$, for an error
of order $1/r^{11}$. The three-term digit counts in Table~\ref{tab:pathology}
and in the certificate are unchanged and correct for that truncation; the
four-term column and the four-term crossover are new. Through v0.16.0 this
paragraph also placed the floor ``beyond a few tens of cells''.)}

\paragraph{The enclosures stay tight where double precision has nothing.}
The point of carrying exact rationals is that the cancellation is defeated, not
merely diagnosed. At the canonical far point, the cube on-axis $N_{xx}$ at
$n=10^{4}$, the certified enclosure has width $\approx 3.6\times10^{-67}$ about a
value of magnitude $1.6\times10^{-13}$ --- a relative width near $2\times10^{-54}$,
i.e. the entry is pinned to some fifty-four correct digits exactly where the
double-precision analytic formula has none. Across the whole table the relative
enclosure width never exceeds $1.7\times10^{-46}$, attained at the elongated cell
on the body diagonal, still some forty-five correct digits. We had pre-registered a kill
condition: if the certified enclosures could not be made tighter than double
precision anywhere in the regime of interest, the target dies. They are tighter
than double precision \emph{everywhere} in the regime, and by tens of digits
where it matters most. The target is \textsc{live}.

\paragraph{Self-demag factors.}
As a byproduct the certificate carries the diagonal self-terms $\mathbf r=0$ ---
the single most-used demag numbers, the thin-film and cube demagnetizing factors
--- as certified rationals: the cube's $N_{xx}(0)=\tfrac13$ exactly, and the film
and rod factors to the full working precision. These have no cancellation
pathology; we include them because a certified reference table of demag
constants is exactly the kind of common, load-bearing, unowned object this
program targets.

\section{What is certified, what is enclosed, what is trusted}\label{sec:trusted}

\paragraph{Machine-checked, replayable with standard-library Python.}
The checker reads only the certificate. It (i)~recomputes the SHA-256 of the
canonical entry list; (ii)~for every entry, recomputes the enclosure by interval
arithmetic and Newell evaluation at the certificate's stated precision --- code
that is not independent of the private engine; see below --- and verifies that
the recomputed rigorous interval lies inside the
certified $[N_{\mathrm{lo}},N_{\mathrm{hi}}]$ --- which pins
$N_{\mathrm{lo}}\le N_{\mathrm{true}}\le N_{\mathrm{hi}}$ --- while rejecting an
enclosure that has been narrowed off the truth, shifted, or implausibly widened;
(iii)~recomputes the naive double value and requires agreement to a relative
$10^{-4}$ (bit for bit with \texttt{--strict-naive}; the low bits of a naive value
depend on the platform's \texttt{libm}), and recomputes each rigorous digit-loss
bracket from the recorded value. \emph{(Correction, v0.17.0, 2026-11-10: through
v0.16.0 clause (iii) said bit for bit; since v0.16.0 that is the
\texttt{--strict-naive} mode.)} Any mismatch prints \texttt{CHECK FAIL} and exits
nonzero. The checker
imports only \texttt{sys}, \texttt{json}, \texttt{math}, \texttt{hashlib},
\texttt{fractions} and \texttt{collections}; it uses no signals, subprocesses,
network or wall-clock, and
runs unchanged under a clean environment.\footnote{Erratum, 2026-08-12: v0.12.0 and v0.12.1 gave this import list without
\texttt{collections}, which \texttt{check\_demag.py} imports for the
consistency-identity pass. The list is now complete; every import remains in the
Python standard library, and nothing about the checker changed.} A tamper
battery --- an enclosure
narrowed to exclude the true value, an enclosure shifted off it, an enclosure
grossly widened, a corrupted naive value, an altered digit-loss claim, a
falsified hash --- is rejected nonzero in every case, at the hash layer and, when
the hash is recomputed to match, at the mathematical re-derivation layer.

\paragraph{Trusted.} Three things, stated plainly.
\begin{enumerate}
  \item \emph{The Newell formulas are the object.} The certificate bounds the
  value of the analytic Newell tensor entry \eqref{eq:newellf}--\eqref{eq:stencil}.
  Whether that analytic entry is the right physical kernel for a given
  discretization --- the energy-based averaging, the rectangular-cell assumption
  --- is Newell's modelling choice \cite{Newell1993}, not something a certificate
  can establish. We bound the number the simulators actually compute.
  \item \emph{CPython, the operating system, and the hardware}, for the checker's
  own exact-rational arithmetic.
  \item \emph{SHA-256 collision resistance}, for the bookkeeping only; a skeptic
  can regenerate and re-check every enclosure without trusting any hash.
\end{enumerate}
The generator that emitted the certificate is not trusted for any value: the
checker recomputes every enclosure from the certificate's own parameters and
rejects any entry it cannot reproduce. It is not, however, code-independent
of the private files that built the certificate. Counted on the v0.15.0 file,
212 of the 558 executable lines of \texttt{check\_demag.py} appear verbatim in
\texttt{gen\_demag.py}, \texttt{civ.py}, \texttt{newell.py} and
\texttt{validate.py}. The naive-double Newell block, lines 393 to 509 of the
shipped file, shares 80 of its 103 lines with
those files, 76 of them with \texttt{newell.py}, and contains byte-identical
runs of 31 and 26 lines against \texttt{newell.py}; the interval Newell
evaluation shares 43 of 143 lines and the certified interval class 56 of 152
with the same files; the digit-loss bracket shares 10 of 28 with
\texttt{gen\_demag.py}. A passing check may therefore not be read, on its
own, as an independent re-derivation of Newell's formulas, of the interval
arithmetic, or of the double-precision pathology. \emph{(Correction, v0.17.0,
2026-11-10: the counts in this section were measured on the v0.15.0 file. The
v0.16.0 portability change brought it to 568 executable lines, and the line
numbers, which through v0.16.0 were those of v0.15.0, are now those of the
shipped file.)}

Two things stand outside that shared lineage, and they are what the result
rests on. The first is external. \texttt{demag-scripts/anchor\_check.py}
takes 16 rows of the 31-row Maple 50-digit check table published in OOMMF's
\texttt{demagcoef.cc} --- a gold standard computed by other people, in other
software, for other purposes --- and requires the checker's own recomputed
enclosure midpoint to agree with each to at least 48 digits. That anchors 16
rows of 862. It is a spot check against an external standard, not coverage
of the table, and the remaining 846 entries have no external anchor. The
second is internal and covers every entry. At lines 663 to 691 the checker
imposes two consistency identities on the certified endpoints themselves,
with no re-derivation involved: the three diagonal components must sum to an
interval enclosing 0 at each of the 136 mutual points, and enclosing 1 at
each of the 4 self-terms. Apart from a bare \texttt{else:}, those lines
share no text with the private files, and neither do the comparisons against
the certificate or the fail-and-exit reporting.

One part of the bookkeeping is shared, necessarily. The canonical
serialisation of an entry and the SHA-256 recomputation, lines 545 to 549
and 580 to 581, are common with \texttt{gen\_demag.py}: both sides must
serialise the entry list identically or no hash could ever agree. The hash
check therefore establishes that the shipped entry list is the one that was
hashed, not that two independent implementations agree on how to hash it.
The tamper battery rejects at the hash layer and, when the hash is made to
match, at the mathematical layer; both remain in force.

\section{The reference table, and certification}\label{sec:cert}

The permanent record is a single certificate JSON, holding the method parameters
(the working precision, the transcendental enclosure scheme, the stencil) and, for
each grid point, the certified enclosure, the naive double value, the value of
the three-term asymptotic truncation where defined (stored under the field name
\texttt{asymp\_oommf}, which predates the correction in \S\ref{sec:results}; the
certificate's \texttt{formulas} and \texttt{method} texts likewise predate the
corrections in \S\ref{sec:object} and \S\ref{sec:results}), and the rigorous
digit-loss brackets. With it goes the standard-library checker, which re-derives
every enclosure, naive value and digit-loss bracket. The checker does not
recompute the stored asymptotic value. \texttt{demag-scripts/four\_term\_digits.py}
recomputes the three-term truncation from the series~\eqref{eq:series} and prints
its correct digits; it does not read the stored value. The two together are the
public unit; a reader with nothing but CPython replays any enclosure from the
certificate alone. The novelty sweep (dated, shipped with the artifacts) found
no certified or interval-arithmetic demagnetization-tensor table in the
literature: the accuracy studies \cite{CF2015,Bjork2023} use floating-point
reference values inside error plots, and the newer codes inherit the same
uncertified kernels. The generator and the search that chose the grid remain
private, per this program's boundary; they are not needed to check anything.
\emph{(Correction, v0.17.0, 2026-11-10: through v0.16.0 this paragraph said the
checker re-derives all of the record, the asymptotic value included.)}

\section{Questions}

\begin{question}[The certified crossover, shipped to the codes]
The certificate brackets, per aspect ratio and direction, the separation at
which a truncated asymptotic expansion overtakes the naive analytic formula. Can this
crossover map be turned into a certified switching rule that a simulator adopts
directly --- a table of ``use analytic below $n^\ast$, asymptotic above'' with
$n^\ast$ rigorously bracketed --- and does the double-double arithmetic OOMMF uses
below the cutoff move $n^\ast$ by exactly the amount the $10^{-15}\to10^{-32}$
precision change predicts? OOMMF already sets its switching radii from
assumed error and speed models (\S\ref{sec:bound}); the question is whether radii
with certified error bounds can be supplied to such a schedule. \emph{(Correction, v0.17.0, 2026-11-10: ``OOMMF's
asymptotic expansion'' replaced; see \S\ref{sec:results}.)}
\end{question}

\begin{question}[Certifying the asymptotic tail]
The asymptotic expansion is itself only enclosed here through its floating-point
value. Can its truncation error be enclosed rigorously, by a certified remainder
bound, so that, together with the analytic enclosure, one obtains a single
certified kernel that is tight at \emph{every} separation, short and long, with a
proof-carrying handover in between? \emph{(Partly answered in
\S\ref{sec:bound}, v0.17.0, 2026-11-10: an explicit bound on the truncation error for every
rectangular cell, tensor component and separation larger than the cell diagonal. The
single certified kernel with a proof-carrying handover remains open. Through
v0.16.0 this question spoke of ``OOMMF's order-$1/r^{6}$ expansion''; see
\S\ref{sec:results}.)}
\end{question}

\begin{question}[Beyond the rectangular prism]
The tetrahedral and cylindrical tiles of \cite{Bjork2023}, and the periodic and
spatially-adaptive kernels of the newer codes, share the same cancellation
structure. Do their governing functions admit the same outward-rounded interval
treatment, and is the breakdown radius a computable function of the tile geometry?
\end{question}

\section{Addendum (v0.17.0, November 2026): a bound on the asymptotic truncation error}\label{sec:bound}

\paragraph{The series.}
Take two identical cells with edges $(\Delta x,\Delta y,\Delta z)$ parallel to the
axes, volume $V=\Delta x\,\Delta y\,\Delta z$ (written $|\sigma|$
in~\eqref{eq:tensordef}) and centre offset $\mathbf r$, with length $r=|\mathbf r|$
and direction $\hat{\mathbf r}$, and let $D=\sqrt{\Delta x^2+\Delta y^2+\Delta z^2}$ be the cell diagonal. In this
section $n$ denotes a truncation order; in \S\ref{sec:results} it counts cells.
Let $\mathbf u$ be the difference of two independent uniform points of
the cell. Each coordinate of $\mathbf u$, for example $u_x$, has the triangular law on
$[-\Delta x,\Delta x]$, so $\mathbb E[u_x^{2m}]=\mu_m\,\Delta x^{2m}$ with
$\mu_m=2/\big((2m+1)(2m+2)\big)$, and likewise for $y$ and $z$. For cells that do
not overlap, \eqref{eq:tensordef} is
\begin{equation}
  N_{ij}(\mathbf r) = -\frac{V}{4\pi}\,
  \mathbb E\Big[\partial_i\partial_j\frac{1}{|\mathbf r+\mathbf u|}\Big],
  \label{eq:repr}
\end{equation}
and for $r>D$, which already rules out overlap, expanding in $\mathbf u$ and
averaging term by term gives
\begin{equation}
  N_{ij}(\mathbf r) = -\frac{V}{4\pi}\sum_{m\ge 0}
  \frac{T_{2m+3}(\hat{\mathbf r})}{r^{2m+3}},
  \label{eq:series}
\end{equation}
where $T_{2m+3}$ is a polynomial in the direction cosines whose coefficients are
polynomials in the squared edges. \texttt{demag-scripts/series\_terms.py}
computes these terms to any order in exact rational arithmetic. The series is not
new. Donahue gives it to all orders in operator form \cite{Donahue2007}, and
OOMMF's \texttt{demagcoef.cc} carries its first four terms, $T_3$ to $T_9$; our
exact derivation reproduces OOMMF's integer coefficients for $N_{xx}$ and
$N_{xy}$. Write $S_n$ for~\eqref{eq:series} truncated after the terms through
$1/r^{n-2}$, so that the omitted terms start at $1/r^{n}$. OOMMF's asymptotic
order is this $n$. $S_n$ is the exact truncated series;~\eqref{eq:boundfree}
and~\eqref{eq:boundshape} below bound its truncation error only, and a
floating-point evaluation adds rounding error.

\begin{theorem}\label{thm:bound}
For all edges $(\Delta x,\Delta y,\Delta z)$, every component $(i,j)$, every odd
$n\ge 3$ and every $r>D$,
\begin{equation}
  \big|N_{ij}(\mathbf r)-S_n(\mathbf r)\big|
  \le \frac{V\,D^{\,n-3}}{2\pi\, r^{\,n-2}\,(r^2-D^2)},
  \label{eq:boundfree}
\end{equation}
and, using the cell shape,
\begin{equation}
  \big|N_{ij}(\mathbf r)-S_n(\mathbf r)\big|
  \le \frac{V}{4\pi}\left[\frac{(n-2)(n-1)\,\mathbb E|\mathbf u|^{n-3}}{r^{n}}
  + \frac{2\,D^{\,n-1}}{r^{n}\,(r^2-D^2)}\right],
  \label{eq:boundshape}
\end{equation}
whose right-hand side never exceeds that of~\eqref{eq:boundfree}. The constant
in~\eqref{eq:boundfree} is attained, for the diagonal component along the axis,
in the limit of a thin needle with offset along its axis, so it cannot be lowered
without using the cell shape.
\end{theorem}

\begin{proof}
Fix $\mathbf u$; then $|\mathbf u|\le D<r$. Write
$\varphi=\partial_i\partial_j(1/|\mathbf x|)$, and let $\mathbf e_i,\mathbf e_j$ be
the coordinate unit vectors. The function $t\mapsto\varphi(\mathbf r+t\mathbf u)$
extends holomorphically to the disc $|t|<r/|\mathbf u|$, whose radius exceeds $1$,
because the zeros of the quadratic $|\mathbf r+t\mathbf u|^2$ in $t$ have modulus
$r/|\mathbf u|$; so its Taylor series in $t$ converges at $t=1$. Its $t^k$
coefficient is the $k$-th term of the Taylor series of $\varphi$ about
$\mathbf r$, evaluated at $\mathbf r+\mathbf u$:
$\tfrac1{k!}\,\nabla^{k+2}(1/|\mathbf x|)(\mathbf r)[\mathbf u,\dots,\mathbf u,
\mathbf e_i,\mathbf e_j]$. The $(k+2)$-th derivative of $1/|\mathbf x|$ is a
symmetric multilinear form on Euclidean space, and the norm of such a form equals
its largest absolute value on the diagonal, over unit vectors
(\cite[Satz~I]{Banach1938}; in finite dimensions it also follows, by induction on
the degree, from the gradient bound of \cite[Theorem~IV]{Kellogg1928}). For a
unit vector $\mathbf v$ at angle $\gamma$ to $\mathbf x$, the Legendre generating
function gives
$(\mathbf v\cdot\nabla)^{\ell}(1/|\mathbf x|)=(-1)^{\ell}\,\ell!\,P_{\ell}(\cos\gamma)/|\mathbf x|^{\ell+1}$,
where $P_{\ell}$ is the Legendre polynomial of degree $\ell$, and
$|P_{\ell}|\le 1$ with equality at $\gamma=0$. The norm of the $\ell$-th derivative
at $\mathbf r$ is therefore exactly $\ell!/r^{\ell+1}$; with $\ell=k+2$, the
$k$-th term is at most $(k+1)(k+2)\,|\mathbf u|^k/r^{k+3}$.
These bounds form a convergent series, so the expansion may be averaged term by
term. The odd terms average to zero because $\mathbf u$ and $-\mathbf u$ have the
same law. For even $k=2s$, $\mathbb E|\mathbf u|^k$ is the multinomial sum of
$\mathbb E[u_x^{2a}]\,\mathbb E[u_y^{2b}]\,\mathbb E[u_z^{2c}]$ over $a+b+c=s$.
One has $\mu_a\mu_b\le\mu_{a+b}$: putting
$A=2a+1$ and $B=2b+1$, the difference $A(A+1)B(B+1)-2(A+B-1)(A+B)$ is a
quadratic in $A$ with nonnegative leading coefficient $B(B+1)-2$, it vanishes at
$A=1$, and its derivative there is $(3B+2)(B-1)\ge 0$. Hence
$\mathbb E|\mathbf u|^{2s}\le\mu_s D^{2s}$, that is
$(k+1)(k+2)\,\mathbb E|\mathbf u|^k\le 2D^k$, with equality for a needle. The
omitted terms of $S_n$ start at $k=n-3$, so
$|N_{ij}-S_n|\le\frac{V}{4\pi}\sum 2D^k/r^{k+3}$ over even $k\ge n-3$, a
geometric series in $D^2/r^2$ whose sum is~\eqref{eq:boundfree}. Keeping the
exact moment for $k=n-3$ gives~\eqref{eq:boundshape}. In the limit of a needle
of vanishing width, with offset and component along its axis, every inequality
above becomes an equality; after division by $V$ both sides
of~\eqref{eq:boundfree} depend continuously on the edges, so the ratio of error
to bound tends to $1$, which gives the last statement.
\end{proof}

\paragraph{Tests.}
\texttt{demag-scripts/check\_remainder\_bound.py} tests~\eqref{eq:boundfree}
and~\eqref{eq:boundshape} at $n=5,7,9,11$ against every certified entry with
$r>D$ in both certificates ($1{,}162$ entries of
\texttt{demag\_certificate\_six\_cells.json}, which include the $776$ of
\texttt{demag\_certificate.json}), in exact rational
arithmetic, with $\pi$ and $r$ bracketed on the safe side and the true value
anywhere in its certified enclosure. No comparison fails. The largest ratio of
error to bound is $0.833$, for~\eqref{eq:boundfree} at $n=5$, the elongated cell
along its axis at $10{,}000$ cells. The script also checks the moment inequality
for every cell. These checks test Theorem~\ref{thm:bound}; the proof is above.
Separately, a test with the private interval engine, which is not
replayable from this repository, covered $2{,}759$ random cases (all six
components, cell aspect ratios up to about $1{,}000$, $r/D$ from $1.0001$ to
$3{,}150$) and found no violation of~\eqref{eq:boundfree}
or~\eqref{eq:boundshape}; the largest ratio was $0.99997$, for a
$1:1/300:1/300$ needle along its axis. The second certificate,
\texttt{demag\_certificate\_six\_cells.json}, extends the first by two thin-film
cells, $(1,1,1/10)$ and $(1,1.3,1/10)$; its $862$ entries for the original four
cells are identical to those of the first, and it passes
\texttt{check\_demag.py} unchanged.

\paragraph{Relation to OOMMF's switching radii.}
Since version 2.0a2, OOMMF's \texttt{FindKnotSchedule} chooses its switching
radii from assumed error and speed models \cite{OOMMF,DP2019}. The one for the base,
unsubdivided asymptotic method is stated in a source comment, in our notation, as
$V\,h_{\max}^{\,n-3}/\big(5\,r^{n-2}(r^2-h_{\max}^2)\big)$, with $h_{\max}$ the
longest edge. Bound~\eqref{eq:boundfree} has the same form, with $D$ in place of
$h_{\max}$ and $1/(2\pi)$ in place of $1/5$. It is loose for compact cells. For a
cube at $n=11$, in units of $(V/4\pi)\,h_{\max}^{8}/r^{11}$, the leading error
coefficient is $99/64$ ($N_{xx}$ with offset along $x$, or any diagonal component
along its own axis; the largest over directions and components); OOMMF's model gives $4\pi/5\approx2.5$, bound~\eqref{eq:boundfree}
gives $162$ and bound~\eqref{eq:boundshape} gives $1122/35$, about $105$ and $21$
times the true coefficient.

\paragraph{Prior work.}
The series is Donahue's \cite{Donahue2007}, and its first four terms are in
OOMMF. The argument above is the three-dimensional form of the classical
multipole truncation estimate, which \cite{GR1987} gives in two dimensions,
applied to the second derivatives of the kernel averaged over the pair of cells. Buchheit, Busse, Ke\ss ler and Rybakov give a rigorous remainder estimate
for the cuboid moment series by Cauchy estimates
\cite[Corollary~2.5]{BBKR2026}; its constant is not explicit and it requires a
separation larger than a multiple of the cell size. A search on 2026-11-10
(record: \texttt{SWEEP-RECORD-DEMAG-BOUND-2026-11-10.md}) found no bound of the
explicit form~\eqref{eq:boundfree}. The search does not rule out a bound of this
form elsewhere.

\begin{thebibliography}{99}
\bibitem{Banach1938} S. Banach, \emph{\"Uber homogene Polynome in $(L^2)$},
Studia Math. \textbf{7} (1938) 36--44.
\bibitem{BBKR2026} A. A. Buchheit, J. K. Busse, T. Ke\ss ler and F. N. Rybakov,
\emph{Zeta expansion for long-range interactions under periodic boundary
conditions with applications to micromagnetics}, J. Comput. Phys. \textbf{559}
(2026) 114885; arXiv:2509.26274.
\bibitem{Bjork2023} R. Bj\o rk and M. d'Aquino, \emph{Accuracy of the analytical
demagnetization tensor for various geometries}, J. Magn. Magn. Mater.
\textbf{587} (2023) 171245; doi:10.1016/j.jmmm.2023.171245.
\bibitem{CF2015} D. S. Chernyshenko and H. Fangohr, \emph{Computing the
demagnetizing tensor for finite difference micromagnetic simulations via
numerical integration}, J. Magn. Magn. Mater. \textbf{381} (2015) 440--445;
arXiv:1403.1978.
\bibitem{demagderiv} K. M. Lebecki, \emph{demagderiv}, Mathematica notebook
(2012) that transcribes the asymptotic demagnetization tensor of a 2010 OOMMF
pre-release (oommf12a4pre-20100909) from its code and differentiates it for the
MFM signal; \texttt{doc/demagderiv.pdf} in R. Boardman, K. M. Lebecki and
H. Fangohr, \emph{ovf2mfm}, \texttt{bitbucket.org/fangohr/ovf2mfm}, commit
494b363f7edb.
\bibitem{Donahue2007} M. J. Donahue, \emph{Accurate computation of the
demagnetization tensor}, talk at HMM-2007, Naples, 4 June 2007; slides at
\texttt{math.nist.gov/\~{}MDonahue}.
\bibitem{DP2019} M. J. Donahue and D. G. Porter, \emph{Quantitative evaluation
and reduction of error in computation of the demagnetization tensor}, poster
ET-01, MMM 2019, Las Vegas, 7 November 2019;
\texttt{math.nist.gov/\~{}MDonahue/talks/mmm2019-poster.pdf}.
\bibitem{GR1987} L. Greengard and V. Rokhlin, \emph{A fast algorithm for particle
simulations}, J. Comput. Phys. \textbf{73} (1987) 325--348;
doi:10.1016/0021-9991(87)90140-9.
\bibitem{Kellogg1928} O. D. Kellogg, \emph{On bounded polynomials in several
variables}, Math. Z. \textbf{27} (1928) 55--64.
\bibitem{Newell1993} A. J. Newell, W. Williams and D. J. Dunlop, \emph{A
generalization of the demagnetizing tensor for nonuniform magnetization}, J.
Geophys. Res. Solid Earth \textbf{98}(B6) (1993) 9551--9555.
\bibitem{OOMMF} M. J. Donahue and D. G. Porter, \emph{OOMMF User's Guide,
Version 1.0}, Interagency Report NISTIR 6376, NIST (1999); source files
\texttt{demagcoef.cc}, \texttt{demagcoef.h}, \texttt{demag.cc} and
\texttt{fft3v.h} in \texttt{app/oxs/ext/} and the platform files
\texttt{config/platforms/*.tcl}, releases 1.2a5bis to 2.1a0 as mirrored at
\texttt{github.com/fangohr/oommf}.
\end{thebibliography}

\end{document}
